% Round 15. New candidate local interface, independent of the curve-base repair.
% Requires independent review. Does not itself assert a Green operator theorem.
\section{Local separable perturbation of rational moments}
\label{r15b:local-perturbation-section}

This section addresses a specific obstruction in extending the
curve-base construction. A fixed global rational tuple can have a
discriminant divisor passing through the chosen base point. The
tuple representing a PWS form is not fixed: sufficiently small
changes of its rational functions preserve their norms on the
compact domains on which they are used. The statements below make
that freedom simultaneous with exact domain membership.

Throughout $K/\mathbf Q_p$ is the original valued field. All
analytic neighborhoods are ordinary neighborhoods. A polynomial
perturbation changes a presenting moment, not the actual form it
represents. No preservation of closedness by an ambient extension
is assumed.

\begin{lemma}[Threshold stability on the whole ambient box]
\label{r15b:threshold-stability}
Let $S$ be any set of valued points, let $u,\widetilde u$ be
functions whose values lie in their corresponding valued residue
fields, and let $r>0$. Suppose
\[
                 |u-\widetilde u|<r\quad\hbox{at every point of }S.
\tag{SP15.1}\label{r15b:threshold-error}
\]
Then each of the four conditions
$|u|<r$, $|u|\leq r$, $|u|>r$, $|u|\geq r$ is equivalent
to the corresponding condition for $\widetilde u$, point by point
on $S$.
\end{lemma}

\begin{proof}
If $|u|\geq r$, the strict inequality in
\eqref{r15b:threshold-error} gives
$|\widetilde u|=|u|$. If $|u|<r$, the ultrametric inequality
gives $|\widetilde u|<r$. The same reasoning with $u$ and
$\widetilde u$ interchanged proves all four equivalences,
including the boundary $|u|=r$.
\end{proof}

\begin{lemma}[Small original-field balls are Zariski dense]
\label{r15b:small-ball-density}
Let $L/K$ be any field extension, and let $B_1,\ldots,B_m$ be
nonempty open valuation balls in $K$. If
$H\in L[T_1,\ldots,T_m]$ is nonzero, then $H$ does not vanish
at every point of $B_1\times\cdots\times B_m$.
In particular, a finite list of nonzero polynomials over $L$
can simultaneously be avoided by arbitrarily small choices of
parameters in the original field $K$.
\end{lemma}

\begin{proof}
Every such ball contains arbitrarily many distinct elements of
$K$, since the valuation is nontrivial. If $d_j$ bounds the degree
of $H$ in $T_j$, choose $d_j+1$ distinct elements of $B_j$.
Vanishing on their Cartesian product would force every coefficient
of $H$ to vanish by successive inversion of the Vandermonde matrices.
Those matrices have nonzero determinants in $K$ and are therefore
invertible over $L$. This contradicts $H\ne0$. For a finite list
apply the result to its nonzero product.
\end{proof}

Let $V=\operatorname{Spec}R$ be a smooth integral affine
$K$-variety, let $x\in V^{\mathrm{an}}$, and put
$p:V\times\mathbf P^1\to V$.
Use the two standard affine coordinates $t_0=z$ and
$t_\infty=w=1/z$ on the projective line.

\begin{definition}[A finite Laurent presentation around a complete fiber]
\label{r15b:Laurent-fiber-presentation}
Such a presentation consists of the following data.
\begin{enumerate}
\item A compact algebraic Laurent affinoid domain
$A\subset V^{\mathrm{an}}$ with $x\in A^\circ$.
Here ``algebraic Laurent'' means a domain given, inside a bounded
affine-coordinate polydisc, by finitely many scalar upper and lower
absolute-value bounds on elements of $R$, with all thresholds
positive.
\item Finitely many boxes
\[
 K_i=A\times\{|t_{\sigma_i}|\leq R_i\},\qquad
        \sigma_i\in\{0,\infty\},\quad R_i>0,
\]
and compact domains $A_i\subset K_i$ defined by finitely many
conditions
\[
      |P_{i\ell}(t_{\sigma_i})|\leq r_{i\ell}
      \quad\hbox{or}\quad
      |P_{i\ell}(t_{\sigma_i})|\geq r_{i\ell},
 \qquad P_{i\ell}\in R[t_{\sigma_i}],\quad r_{i\ell}>0.
\tag{SP15.2}\label{r15b:Laurent-inequalities}
\]
The ordinary interiors of the $A_i$ cover the entire analytic
fiber $p^{-1}(x)$.
\item Each of a specified finite family of PWS forms on
$(V\times\mathbf P^1)^{\mathrm{an}}$ has a presentation on an
ordinary neighborhood of $A_i$ in finitely many analytic units.
Each of those units has been replaced, with exactly the same norm
on an ordinary neighborhood of $A_i$, by a rational function
\[
          R_{ik}=N_{ik}(t_{\sigma_i})/Q_{ik}(t_{\sigma_i}),
       \qquad N_{ik},Q_{ik}\in R[t_{\sigma_i}],
\tag{SP15.3}\label{r15b:rational-unit}
\]
where both $N_{ik}$ and $Q_{ik}$ are nonzero at every point of
$A_i$.
\end{enumerate}
Repeated occurrences of the same polynomial are allowed to be
listed separately. The fixed box coordinate and its bound are
not among the polynomials to be perturbed.
\end{definition}

\begin{lemma}[Existence of a finite Laurent fiber presentation]
\label{r15b:Laurent-fiber-exists}
Given $V,x$ and a finite family of ordinary PWS forms on
$(V\times\mathbf P^1)^{\mathrm{an}}$, a presentation as in
Definition~\ref{r15b:Laurent-fiber-presentation} exists. There is
an ordinary neighborhood $U$ of $x$, with $U\subset A^\circ$,
such that
\[
                 p^{-1}(U)\subseteq\bigcup_i A_i^\circ.
\tag{SP15.4}\label{r15b:whole-fiber-tube}
\]
The same assertion holds if the forms are initially given only
over the inverse image of an ordinary base neighborhood of $x$.
\end{lemma}

\begin{proof}
First choose $A$ inside the prescribed base neighborhood. At each
point of the compact fiber $p^{-1}(x)$ choose one of the two affine
fiber charts and a simultaneous ordinary presentation domain for
the finitely many forms. The topology on an affine algebraic
analytification is generated by absolute values of regular
functions. Finitely many strict scalar inequalities therefore give
an ordinary neighborhood inside that presentation domain.
Tighten their positive bounds while retaining the chosen point in
the ordinary interior, and add a sufficiently large fixed bound
on the affine fiber coordinate. Intersect with the inverse image
of $A$. This gives a domain $A_i$ of the stated form inside the
presentation domain. Compactness of the fiber gives finitely many
such $A_i$.

The coordinate algebra of an algebraic Laurent domain is obtained
by completing the regular coordinate algebra after adjoining the
inverses of the finitely many functions given positive lower
bounds. Thus every analytic unit on $A_i$ can be approximated
uniformly by rational functions with denominators that are products
of those functions. The base-domain lower-bound functions are
included among these denominators. Choose the approximation error
strictly below the positive minimum norm of the unit on $A_i$.
The denominator is nonzero on $A_i$, and the ultrametric equality
makes the numerator nonzero there as well. This gives
\eqref{r15b:rational-unit}. Both the original unit and the rational
function are analytic on an ordinary neighborhood of $A_i$;
the strict error inequality persists on such a neighborhood.
Their norm functions, and therefore all their PWS moment
pullbacks, agree there. This is also the frozen R14 exact-norm
replacement argument in this explicit Laurent situation.

The union $N=\bigcup_i A_i^\circ$ is an ordinary open containing
the entire fiber. The projection $p$ is proper, so
$p((V\times\mathbf P^1)^{\mathrm{an}}\setminus N)$ is closed
and does not contain $x$. Intersect its open complement with
$A^\circ$ to obtain $U$. This proves
\eqref{r15b:whole-fiber-tube} without assuming that the topology of
an analytic fiber product is the Cartesian-product topology.
\end{proof}

\begin{proposition}[Simultaneous norm preservation and root separation]
\label{r15b:separable-perturbation}
Given a finite Laurent fiber presentation, every polynomial
occurrence in \eqref{r15b:Laurent-inequalities} and
\eqref{r15b:rational-unit} can be replaced by a polynomial
$\widetilde P\in R[t_{\sigma_i}]$ by adding arbitrarily small
coefficients from $K$, with the following properties.
\begin{enumerate}
\item Replacing the domain polynomials leaves each $A_i$
exactly unchanged as a subset of $K_i$, hence as a domain over
$A$. This equality holds for points outside $A_i$ as well as
for points inside it.
\item For every presenting rational unit,
\[
       \left|\frac{\widetilde N_{ik}}{\widetilde Q_{ik}}\right|
                   =|R_{ik}|
\]
on an ordinary neighborhood of $A_i$. Therefore the original
actual PWS forms have the same local tensor presentations after
the replacements.
\item Set $D_P\geq\max(1,\deg P)$ for each occurrence before
choosing the perturbation, and allow all coefficients through
degree $D_P$ to vary. Express the perturbed polynomials in the
single coordinate $z$ by
\[
 H_P(z)=\widetilde P(z)\quad(\sigma_i=0),\qquad
 H_P(z)=z^{D_P}\widetilde P(1/z)\quad(\sigma_i=\infty).
\tag{SP15.5}\label{r15b:single-coordinate-polynomial}
\]
At the chosen, possibly nonrigid, base point $x$, every $H_P$
has degree $D_P$, has nonzero constant term, and the product
of all the $H_P$ is squarefree over $\mathcal H(x)$.
Thus all their roots are simple, mutually distinct, and
different from $0$ and $\infty$ after splitting.
\end{enumerate}
No algebraic relations between distinct polynomial occurrences
are required to survive the perturbation.
\end{proposition}

\begin{proof}
For one polynomial occurrence on the box $K_i$, a perturbation
$\delta P(t)=\sum_{j=0}^{D_P}\delta_jt^j$ satisfies
\[
        \sup_{K_i}|\delta P|
                 \leq\max_{0\leq j\leq D_P}|\delta_j|R_i^j.
\tag{SP15.6}\label{r15b:uniform-box-bound}
\]
For a domain polynomial require this bound to be strictly below
every threshold at which that occurrence is used. For a numerator
or denominator require it to be strictly below its positive
minimum norm on the corresponding $A_i$. These are finitely
many strict positive bounds on finitely many parameters from $K$,
so they define a product of nonempty arbitrarily small open balls.

The threshold-stability lemma proves each inequality in
\eqref{r15b:Laurent-inequalities} equivalent to its perturbed
inequality on the whole box $K_i$. The fixed box condition itself
is unchanged. Thus the domain is exactly unchanged, including
all of its boundary and all formerly excluded points of the box.
Outside the box both definitions exclude the point. This is why
an estimate only on $A_i$ would not suffice for this part of the
proof.

For presenting unit numerators and denominators, the separate
minimum-norm bounds give
$|\widetilde N_{ik}|=|N_{ik}|$ and
$|\widetilde Q_{ik}|=|Q_{ik}|$ on $A_i$.
They remain units on an ordinary neighborhood of $A_i$, where
the same strict inequalities still hold. This proves the second
assertion. It does not assert equality of the rational functions
themselves, which is unnecessary for PWS moment presentations.

It remains to impose the algebraic generic conditions without
violating any of these smallness bounds. Evaluate all coefficient
functions at $x$, in the field $L=\mathcal H(x)$. Reversal of
the coefficients in the infinity chart is an invertible linear
transformation. Consequently the independent perturbation
parameters allow every coefficient of every polynomial $H_P$
to vary independently over $L$. The product of their leading
coefficients, constant coefficients, and the discriminant of
their total product is a nonzero polynomial in these parameters
with coefficients in $L$; an empty polynomial list is harmless
and requires no choices. For a nonempty list it is nonzero because one can choose,
over the infinite field $L$, monic polynomials of the prescribed
degrees with all roots distinct and nonzero across the entire
finite family. The corresponding coefficient differences give
one parameter value where the displayed polynomial is nonzero.

Lemma~\ref{r15b:small-ball-density} now supplies a parameter value
in the previously specified small balls in the original $K$
where that polynomial is nonzero. All norm and membership
conditions already proved survive. The conditions at $x$ are
exactly the third assertion. In particular the argument is valid
when $x$ has a nonalgebraic completed residue field; it does not
approximate elements of $\mathcal H(x)$ by elements of $K$.
\end{proof}

\begin{corollary}[A splitting cover etale at the chosen base point]
\label{r15b:local-etale-splitting}
After the preceding replacements there is a Zariski open
$V^g\subset V$ with $x\in(V^g)^{\mathrm{an}}$, and a finite
etale surjective map $a:V'\to V^g$ of smooth integral
$K$-varieties, such that every perturbed polynomial splits
into distinct linear factors on $V'$. Their finite roots
$a_1,\ldots,a_m$ are regular units on $V'$, and every nonzero
pairwise difference is a unit. The fixed roots $0$ and
$\infty$ may be adjoined as additional marks.

After shrinking the ordinary base neighborhood $U$ from
\eqref{r15b:whole-fiber-tube} inside $(V^g)^{\mathrm{an}}$,
all these properties hold above its entire inverse image.
No discriminant divisor, root collision, or root pole passes
through this working neighborhood.
\end{corollary}

\begin{proof}
Invert the finitely many leading and constant coefficients and
the discriminant of the total product. They are regular
functions with nonzero values at $x$, so this defines a Zariski
open containing $x$. The normalized product polynomial is monic
with invertible discriminant there. Adjoining a root of a monic
polynomial $F$ with unit discriminant gives a finite free etale
algebra: $F'$ is a unit modulo $F$, the standard etale criterion.
If $F=(T-a)G$ after that adjunction, the identity
$\operatorname{Disc}(F)=\operatorname{Disc}(G)G(a)^2$
shows that $G$ again has unit discriminant. Repeating therefore
constructs a finite etale surjective cover on which all roots
are ordered and split. The standard etale criterion is Stacks
Project, Tag~03PA, Definition~59.26.3 and its following proof.
Choose a connected
component. Because $V^g$ is integral, every nonempty connected
component is still surjective; it is smooth and integral.

The roots are regular because they are roots of the monic
polynomial on that cover. Their product is a unit, as is their
discriminant product, so each root and each pairwise difference
is a unit. The leading coefficients retain the unit factors in
the factorization. The asserted ordinary shrinking follows
because the chosen functions are nonzero at $x$. The cover
may enlarge the constant field; its finite native degree
includes that enlargement. No base modification is involved.
\end{proof}

\begin{proposition}[The local noncompact source image is a finite tree family]
\label{r15b:local-whole-source-presentation}
Retain the preceding data, and let $B$ be a compact affinoid
ordinary neighborhood of $x$ contained in the shrunken $U$.
Put $B'=a^{-1}(B)$, append the fixed finite mark $a_0=0$, and
choose a finite base unit tuple $b$ on $V'$ containing all
leading factors and all nonzero root differences, including the roots
themselves. Set
\[
 P=t_b(B'),\qquad
 W=\bigl(B'\times\mathbf A^1_z\bigr)
                   \setminus\bigcup_{j=0}^m\{z=a_j\}.
\tag{SP15.7}\label{r15b:local-W}
\]
Then the image of the genuine unit tuple
$(b,z-a_0,\ldots,z-a_m)$ on $W$ is precisely the finite
closed polyhedral marked-tree family over $P$ with distance
coordinates $v(a_i-a_j)$. Each of the original forms is
represented on $W$ by one PWS tensor on that finite polyhedral
image, obtained by gluing a fixed finite family of local
expressions. This supplies ordinary actual presentations on the
ordinary source locus over $(B')^\circ$; no statement about a
closed ambient lift or an operator output is made.
\end{proposition}

\begin{proof}
The moment image $P$ is finite polyhedral by affinoid
polyhedrality. At any point of $W$, choose $i$ maximizing
$t=v(z-a_i)$. The ultrametric inequality and maximality give
\[
             v(z-a_j)=\min(t,v(a_i-a_j))
                       \quad(j\ne i).
\tag{SP15.8}\label{r15b:all-point-tree-image}
\]
In the equal-valuation case, strict cancellation would produce
a larger value for $v(z-a_j)$, contradicting the choice of $i$.
Thus every source moment value lies in the displayed tree
family. Conversely, for any actual base point with moment
value $u\in P$ and any real $t$, the Gauss point of radius
$e^{-t}$ centered at $a_i$ in its projective-line fiber has
exactly the values in \eqref{r15b:all-point-tree-image}.
It is a point of $W$. Hence the image is the whole family.
There are finitely many affine comparisons among the supplied
distance coordinates; the family is a finite union of closed
polyhedra. This statement concerns the entire moment image,
not only the presenting skeleton or almost every fiber point.

Every perturbed domain polynomial and presenting rational unit
has valuation affine in the common tuple: factor it into the
split root functions and its base unit factor, and use
\eqref{r15b:single-coordinate-polynomial} in the infinity chart.
The fixed box coordinate is $z$ or $1/z$, and its bound is
also an affine inequality in this tuple. Membership in each
pulled-back $A_i$ is therefore determined by a fixed closed
polyhedron $Q_i$ in the common moment space. The domains cover
the whole source over $B'$ by \eqref{r15b:whole-fiber-tube}.
The moment image of an overlap is exactly the overlap of the
corresponding image pieces.

The local coefficient expressions can first be extended to
globally defined Euclidean PWS tensors in their own variables
without altering their values near the compact moment image
of $A_i$. This is the degreewise finite-polyhedron extension
lemma and a Euclidean smooth cutoff. Only the coefficients
used for that particular domain need be bounded; other common
moment coordinates may be unbounded on its image piece.
After affine substitution the finitely many expressions are
PWS on the common finite polyhedral image. Their pullbacks
agree as the original actual form. The frozen R14 fixed-moment
faithfulness argument makes them tangentially equal on every
overlap cell, so they glue.

One can check this last assertion using only compact domains:
intersect the full common moment image with a large cube.
Its inverse image in $W$ is an affinoid Laurent domain over
$B'$, since all root moments are bounded above and below there;
it is compact. Such cubes exhaust the finite image. The
compact-domain faithfulness proof applies on each cube, with
the same finite coefficient expressions. This also checks
lower-dimensional maximal overlap cells. Thus no extension
of equality from a compact domain to an unspecified ordinary
open is being used.

Finally the degreewise extension near any compact portion of
this image, pulled back through genuine units, gives an
ordinary PWS tensor. Mihatsch Proposition 4.7 identifies it
with an ordinary actual form of delta codimension zero.
No derivatives or closedness conditions were lost: the
presentations represent the original form exactly wherever
they are used.
\end{proof}

\paragraph{Scope of the new interface.}
The statements prove that the local rational-presentation and
root-splitting obstruction can be removed near any chosen analytic
base point, in arbitrary smooth base dimension, by changing
presenting moments within their exact-norm classes. They do not
by themselves assert higher-dimensional relative Green outputs,
a projective-bundle Aeppli formula, or modification descent in
higher dimension. Such conclusions still require an explicit
assembly with the root and diagonal cutoffs, the all-test native
integration argument, and finite descent on the local covering.
In particular, the older fixed-global-tuple analysis of bad
divisors is not incorrect: the local perturbed tuple is different.
